\section{Paper2Agent：把论文从 PDF 变成可调用能力}

第二部分讨论的 Paper2Agent 指向一个更大的问题：科研知识如何传播。Zou 指出，传统论文是``passive artifact''，即便附代码，使用门槛仍高。Paper2Agent 的目标是把方法、代码、数据与使用流程打包成可直接对话调用的 Agent。

图\ref{fig:paper2agent-main} 是该部分主流程图。\footnote{来源：Berkeley RDI 课程视频，关键帧时间戳 00:29:27。}

\begin{figure}[H]
\centering
\includegraphics[width=0.95\textwidth]{frame-08-mcp-pipeline.jpg}
\caption{Paper2Agent 自动化流程：从论文和代码到 MCP，再到 Paper Agent}
\label{fig:paper2agent-main}
\end{figure}

\begin{knowledgebox}{读图：Paper2Agent 的可靠性来自测试闭环}
流程从论文与代码仓库开始，Environment Agent 固定依赖，Extraction Agent 暴露方法接口，Testing Agent 用测试验证，再封装成 MCP server。连接成功只是必要条件；若测试不能复现论文行为，Paper Agent 应拒绝把未经验证的函数包装成可信能力。
\end{knowledgebox}

\subsection{从文献阅读到工具可用的中间断层}

多数科研复现痛点不在``看不懂论文摘要''，而在``环境配不起来、脚本跑不动、参数含义不明确、数据前处理缺失''。Paper2Agent 的价值是把这段最耗时、最易出错的中间层结构化为机器流程。

\begin{knowledgebox}{Paper2Agent 典型流水线}
\begin{itemize}
    \item Environment Agent：重建论文代码环境，固定依赖。
    \item Extraction Agent：抽取核心工具接口与关键脚本。
    \item Testing Agent：验证能否重现实验结论。
    \item MCP Packaging：把可调用资源封装为标准协议接口。
\end{itemize}
\end{knowledgebox}

\subsection{为什么是 MCP，而不是临时脚本拼接}

课程给出非常工程化的答案：MCP 提供了统一的资源暴露方式，使下游聊天 Agent 可以稳定调用论文能力，而不是每次重新写胶水代码。这样``论文复用''从一次性劳动变成可持续基础设施。

\begin{importantbox}{Paper MCP 的三个组成}
\begin{itemize}
    \item 工具层：可执行函数与参数约束。
    \item 资源层：数据、文档、复现脚本、补充材料。
    \item 工作流层：任务顺序、失败恢复、结果摘要模板。
\end{itemize}
\end{importantbox}

\subsection{对研究者工作方式的直接影响}

传统模式下，研究者要先读长 PDF，再下载 repo、搭环境、改配置。Paper2Agent 模式下，研究者可先表达任务目标，由 Paper Agent 反向调度工具并返回结构化结果。这里不是取消阅读，而是把``机械复现步骤''优先自动化，让人把精力转向问题选择与结果解释。

图\ref{fig:paper2agent-discussion} 给出了课程讨论页。\footnote{来源：Berkeley RDI 课程视频，关键帧时间戳 00:36:20。}

\begin{figure}[H]
\centering
\includegraphics[width=0.9\textwidth]{frame-03-paper-to-agent.jpg}
\caption{Paper2Agent 讨论要点：把被动论文转成交互 Agent}
\label{fig:paper2agent-discussion}
\end{figure}

\begin{knowledgebox}{读图：从 passive paper 到 interactive agent}
这页把价值分成三层：读者更容易使用方法与数据，知识传播从文本变成接口，多个论文 Agent 还能组合。最后一层风险最高，因为方法假设、数据语义和许可证边界可能不兼容；协调器必须检查前置条件，而不是看到两个 MCP 都可调用就自动串联。
\end{knowledgebox}

\begin{warningbox}{Paper2Agent 并非``任何论文都能一键变 Agent''}
Zou 明确提到并非所有论文都能可靠转换：文档缺失、代码不可运行、数据访问受限、实验步骤未公开都会导致 MCP 质量下降。Paper2Agent 在某种意义上也充当了``可复现性压力测试器''。
\end{warningbox}

\subsection{本章小结}

Paper2Agent 的突破点是知识载体升级：从可读文本升级为可调用能力。它不是替代论文，而是让论文在实践层面真正``可用''，并用 MCP 把可用性标准化。

